\chapter{Security and Compliance}

This chapter presents the security architecture of the platform, tracing each
implemented control from the threat model through to its realisation in code
and its mapping to the OWASP Top Ten 2021~\cite{owasp2021}. Security was
incorporated as a first-class design concern rather than a post-hoc addition:
each API route, data model, and AI integration boundary was evaluated against
the threat model before implementation, and each non-functional requirement
relating to security (NFR-03 through NFR-06 in Table~\ref{tab:nonfunc_reqs})
drove specific architectural decisions documented below.

% ─────────────────────────────────────────────────────────────────────────────
\section{Threat Model}
\label{sec:threat_model}

The platform processes three categories of sensitive assets: user credentials
(email addresses, hashed passwords, and JWT session tokens); educational
media (uploaded video files, generated caption files, and H5P content
packages); and AI-generated assessment artefacts (LLM-suggested questions
and accepted H5P interaction parameters). Following the structured threat
modelling approach recommended by OWASP~\cite{owasp2021}, the following
threat actors and attack surfaces were identified prior to implementation:

\begin{itemize}
  \item \textbf{Unauthenticated external actors} attempting credential brute
  force against the login endpoint, injection via unauthenticated API routes,
  or application-layer denial of service by flooding the video upload or AI
  analysis endpoints with high-volume requests.

  \item \textbf{Authenticated non-administrative users} attempting to read or
  modify another user's videos, H5P content, or account settings through
  Insecure Direct Object Reference (IDOR)---that is, by substituting a known
  UUID for their own resource identifier in API path parameters.

  \item \textbf{Malicious or schema-violating AI responses}: the AI providers
  (Claude and Ollama) are external services whose output cannot be assumed to
  be syntactically or semantically valid. A response that bypasses validation
  could corrupt the H5P content store, inject executable content via
  question text fields, or cause the export pipeline to produce malformed
  \texttt{.h5p} packages.

  \item \textbf{Insecure integration boundaries} in the LTI~1.3 launch
  endpoint, where a crafted or replayed launch request could bypass
  authentication or impersonate an elevated-privilege user within the LMS
  consumer context.
\end{itemize}

% ─────────────────────────────────────────────────────────────────────────────
\section{Authentication}
\label{sec:auth}

\subsection{Password Storage}

Passwords are stored using the bcrypt adaptive hashing algorithm, as
recommended by OWASP for credential storage in web
applications~\cite{owasp2021}. The implementation uses the
\texttt{bcryptjs} library with a cost factor of~12. Bcrypt's adaptive
design means that the computational cost of each hash operation can be
increased as hardware improves, ensuring that the stored credentials remain
resistant to brute-force attack over time. The algorithm incorporates a
cryptographically random per-user salt, preventing rainbow-table and
batch-lookup attacks against the credential store. Plain-text passwords are
never written to disk, log files, environment variables, or API responses
at any point in the registration or authentication flow. This control
directly satisfies NFR-03 and addresses OWASP~A02 (Cryptographic Failures).

\subsection{JSON Web Tokens}

Upon successful credential verification, the backend issues a JSON Web Token
(JWT) conforming to RFC~7519~\cite{jwt2015}. The token is signed using the
HMAC-SHA256 (HS256) algorithm with a 256-bit secret loaded from the server's
environment variables, satisfying NFR-04. The token payload carries the
authenticated user's UUID (\texttt{sub}), username, email, and role, and is
issued with a seven-day expiry (\texttt{exp} claim). The token is delivered
to the client as a JSON response body field and is subsequently attached to
all protected API requests as a \texttt{Bearer} credential in the
\texttt{Authorization} header.

Token verification is implemented as a reusable Express middleware
(\texttt{authenticateToken}) that intercepts every request to a protected
route, verifies the signature and expiry against the server secret, and
attaches the decoded payload to the request context. Requests with absent,
expired, or tampered tokens receive \texttt{HTTP~401 Unauthorized} without
revealing whether the token was absent or invalid---consistent with the
error disclosure minimisation principle~\cite{owasp2021}. This control
addresses OWASP~A01 (Broken Access Control) and A07 (Identification and
Authentication Failures).

\subsection{Email Verification}

Registration flows do not immediately activate new accounts. Instead, a
six-digit numeric verification code is generated using
\texttt{crypto.randomInt}, which provides cryptographically uniform random
integers, and stored with a ten-minute expiry. The code is delivered to the
registered email address via an automated Nodemailer message. The
registration endpoint returns a success response only after the correct code
is submitted to the verification endpoint within the expiry window; accounts
that fail to verify within this window cannot log in. This mechanism ensures
that the email address provided at registration is valid and under the
registrant's control, preventing account enumeration via disposable addresses
and reducing the risk of credential misuse following registration.

\subsection{Password Reset}

The password reset flow uses a 64-character hexadecimal token generated by
\texttt{crypto.randomBytes(32).toString('hex')}, which provides 256~bits of
entropy from the Node.js CSPRNG. The raw token is emailed to the registered
address as part of a one-time reset link; the stored representation is a
bcrypt hash of the token (using the same cost factor as password storage),
ensuring that database compromise does not expose usable reset credentials.
The token carries a one-hour time-to-live stored in the
\texttt{resetTokenExpiry} column. On submission, the backend verifies the
token against the stored hash, checks the TTL, and immediately clears the
token from the database after a single successful use to prevent replay.
This design follows the stateless token pattern recommended by
OWASP~\cite{owasp2021} for credential recovery flows.

% ─────────────────────────────────────────────────────────────────────────────
\section{Authorisation}
\label{sec:authz}

\subsection{Role-Based Access Control}

The platform implements two roles: \texttt{user} (the default educator role)
and \texttt{admin}. Role assignment is stored in the \texttt{Users} table
and embedded in the JWT payload at issuance time. All routes under the
\texttt{/api/admin} prefix are protected by a second middleware layer
(\texttt{requireAdmin}) that inspects the role claim in the validated token
and returns \texttt{HTTP~403 Forbidden} if the claim does not equal
\texttt{admin}. Role-based access control of this form is identified by
OWASP as the appropriate mitigation for A01 (Broken Access Control) in
multi-user web applications~\cite{owasp2021}.

Educator-scoped mutations (updating or deleting a video or H5P interaction)
are additionally protected by ownership checks at the service layer: the
\texttt{userId} foreign key of the retrieved database record is compared
against the authenticated user's UUID before any write operation is executed.
This prevents IDOR attacks in which a user supplies a valid UUID for a
resource belonging to another account.

\subsection{Admin Console Capabilities}

The administrative console, accessible only to users whose JWT carries the
\texttt{admin} role claim, exposes four categories of privileged operation
that are each individually gated by the \texttt{requireAdmin} middleware:

\begin{itemize}
  \item \textbf{User management:} listing all accounts with their
  registration metadata, last-login timestamp, and active status; activating
  or deactivating accounts via an \texttt{isActive} soft toggle; and
  modifying role assignments.

  \item \textbf{Audit logging:} a chronological, append-only record of all
  admin-initiated mutations, recording the acting administrator's UUID,
  action type (drawn from a closed enum), target user UUID, a JSON snapshot
  of the change, and a wall-clock timestamp.

  \item \textbf{System settings:} runtime-configurable parameters stored in
  the \texttt{SystemSettings} table, including maximum upload size, session
  timeout duration, and maintenance mode. These parameters are modified
  exclusively through the admin console, not through environment variables,
  enabling live reconfiguration without server restart.

  \item \textbf{Login tracking:} per-account recording of last successful
  login timestamp and cumulative failed login count, surfaced in the admin
  console user table to support detection of suspicious credential activity.
\end{itemize}

% ─────────────────────────────────────────────────────────────────────────────
\section{Input Validation and Injection Prevention}
\label{sec:validation}

\subsection{Server-Side Request Validation}

All API request bodies pass through \texttt{express-validator} middleware
before reaching any business logic handler, as recommended for injection
prevention under OWASP~A03~\cite{owasp2021}. Validation rules specify field
presence, data type, length bounds, format constraints (e.g., RFC~5322 email
syntax, UUID~v4 format for resource identifiers), and enum membership for
categorical fields (e.g., video processing status). A validation failure
produces an \texttt{HTTP~422 Unprocessable Entity} response containing a
structured error list; no handler code executes. This ``validate first,
process never'' pattern, described by Sommerville as a foundational
principle of defensive software engineering~\cite{sommerville2016}, ensures
that only well-formed inputs reach the data access layer.

\subsection{AI Output as Untrusted Input}

AI-generated suggestion payloads are treated as untrusted external input
regardless of provider, consistent with the principle that probabilistic
model outputs must be validated before integration with deterministic
systems. Each suggestion array returned by the LLM is passed through
\texttt{JSON.parse} and then validated against a Zod TypeScript schema that
enforces the following invariants:

\begin{itemize}
  \item \texttt{timestamp}: a non-negative finite number (seconds from
  video start).
  \item \texttt{type}: a member of the closed set of six supported H5P
  content type identifiers (see Table~\ref{tab:h5p_types}).
  \item \texttt{config}: a non-null record object conforming to the
  content-type-specific parameter schema.
  \item \texttt{reason}: a non-empty string providing pedagogical
  justification for the placement.
\end{itemize}

Suggestions that fail Zod validation are discarded and logged at the
\texttt{warn} level; the remaining valid suggestions are returned to the
client. This prevents schema-violating AI output from corrupting the H5P
content store and ensures that every accepted suggestion conforms to a
known, verifiable schema before injection into the database---addressing
the AI boundary threat identified in Section~\ref{sec:threat_model}.

\subsection{Parameterised Database Access}

All database interactions are executed through the Sequelize
ORM~\cite{sequelize2023}, which generates parameterised SQL queries
internally. Direct interpolation of user-supplied strings into SQL is
absent from the codebase; raw query execution (\texttt{sequelize.query}
with string concatenation) is not used. This architectural constraint
eliminates SQL injection as an attack vector~\cite{owasp2021}, addressing
OWASP~A03 at the persistence layer.

% ─────────────────────────────────────────────────────────────────────────────
\section{HTTP Security Headers}
\label{sec:headers}

The Helmet middleware~\cite{helmet2023} is mounted before all route handlers
in the Express application, configuring a comprehensive set of HTTP security
headers on every response. Table~\ref{tab:helmet} lists the headers
configured and their security function.

\begin{table}[H]
  \centering
  \caption{HTTP security headers configured via Helmet~\cite{helmet2023}}
  \label{tab:helmet}
  \begin{tabularx}{\linewidth}{lX}
    \toprule
    \textbf{Header} & \textbf{Configuration and purpose} \\
    \midrule
    \texttt{Content-Security-Policy}       & Restricts resource loading to declared origins; mitigates cross-site scripting (XSS) via injected external scripts. \\
    \texttt{Strict-Transport-Security}     & \texttt{max-age=31536000; includeSubDomains} --- instructs browsers to enforce HTTPS for one year, preventing protocol downgrade attacks. \\
    \texttt{X-Frame-Options}               & \texttt{DENY} --- prohibits embedding the application in an \texttt{<iframe>}, preventing clickjacking. \\
    \texttt{X-Content-Type-Options}        & \texttt{nosniff} --- prevents browsers from MIME-type sniffing response bodies, eliminating a class of content-injection attacks. \\
    \texttt{Referrer-Policy}               & \texttt{no-referrer} --- suppresses the \texttt{Referer} header on cross-origin requests, preventing URL leakage to third parties. \\
    \texttt{Permissions-Policy}            & Disables camera, microphone, and geolocation access, limiting the browser API surface available to injected scripts. \\
    \bottomrule
  \end{tabularx}
\end{table}

These headers collectively satisfy NFR-05 and address OWASP~A05 (Security
Misconfiguration). Their configuration follows the recommendations published
by OWASP~\cite{owasp2021} and the Helmet project's own security
guidelines~\cite{helmet2023}.

% ─────────────────────────────────────────────────────────────────────────────
\section{Rate Limiting and Request Size Controls}
\label{sec:rate_limit}

The \texttt{express-rate-limit} middleware limits each IP address to
100~requests within any 15-minute sliding window across all API routes.
Requests that exceed this threshold receive \texttt{HTTP~429 Too Many
Requests} with a \texttt{Retry-After} header indicating when the window
resets. Rate limiting of this form is the primary recommended control for
OWASP~A07 (Identification and Authentication Failures), specifically the
brute-force credential attack vector, in which an automated client
systematically attempts password combinations against the login
endpoint~\cite{owasp2021}.

Request body size is bounded at 10~MB for JSON payloads via the
\texttt{express.json} limit option, and at 500~MB for multipart video
uploads via the Multer middleware configuration. Requests exceeding these
limits are rejected before routing, preventing both memory exhaustion under
sustained upload attack and the application-layer denial-of-service risk
associated with unbounded JSON body parsing~\cite{owasp2021}.

% ─────────────────────────────────────────────────────────────────────────────
\section{CORS Policy}
\label{sec:cors}

Cross-Origin Resource Sharing (CORS) is configured with an explicit
origin allowlist enforced by the \texttt{cors} Express middleware. In the
development environment, only \texttt{http://localhost:3002} (the Vite
development server, as noted in Appendix~A) is permitted. In the production
deployment, only the configured Railway domain is permitted. All
preflight (\texttt{OPTIONS}) requests receive the appropriate
\texttt{Access-Control-Allow-*} headers for permitted origins; requests
from unlisted origins are rejected before the authentication middleware
executes. Credentials---specifically the \texttt{Authorization} header
carrying the JWT---are permitted only from explicitly listed origins.
This policy ensures that the API cannot be accessed by scripts loaded from
third-party pages, addressing a cross-origin data exposure vector~\cite{owasp2021}.

% ─────────────────────────────────────────────────────────────────────────────
\section{Security Logging and Audit Trail}
\label{sec:logging}

Structured application logging is implemented with the Winston library,
which writes JSON-formatted log records to both stdout and a rotating log
file. All HTTP requests are logged with method, path, status code, response
time, and authenticated user UUID (when available). Authentication events---
successful login, failed login attempt, password reset request, and email
verification submission---are recorded at the \texttt{info} level with
actor identification and wall-clock timestamp.

Administrative mutations are recorded in a dedicated \texttt{AuditLog}
database table in addition to the application log. Each record carries the
acting administrator's UUID, action type (drawn from a closed string enum),
target user UUID, a JSON snapshot of the changed state, and timestamp. The
audit table is append-only from the application's perspective: no API route
provides a mechanism to update or delete audit records. This design
supports forensic analysis and provides a baseline for institutional
compliance reviews, as recommended for OWASP~A09 (Security Logging and
Monitoring)~\cite{owasp2021}.

Account deletion is implemented as a soft-delete operation: the
\texttt{isActive} flag is set to \texttt{false} and the account is
excluded from all active-user queries, but the record and all associated
video and interaction metadata are retained for audit continuity.
Permanent deletion requires a separate, explicit administrative action.

% ─────────────────────────────────────────────────────────────────────────────
\section{LTI Integration Security}
\label{sec:lti}

The LTI~1.3 launch endpoint (\texttt{POST /api/lti/launch}) validates the
OAuth~1.0a HMAC-SHA1 signature of every incoming launch request before
processing, following the security model specified by IMS
Global~\cite{lti2019}. The consumer key and shared secret are stored
exclusively as server-side environment variables; no LTI credential is
stored in the database or exposed through any API route. Replay protection
is implemented by checking the \texttt{oauth\_timestamp} (within a five-minute
window) and the \texttt{oauth\_nonce} (verified against a short-term nonce
cache) on each launch request~\cite{lti2019}. Requests that fail signature
verification or replay checks receive \texttt{HTTP~400 Bad Request} without
detailed error disclosure.

% ─────────────────────────────────────────────────────────────────────────────
\section{OWASP Top Ten Compliance Summary}
\label{sec:owasp_summary}

Table~\ref{tab:owasp} maps each relevant OWASP Top Ten 2021 category to the
controls implemented in the platform, with cross-references to the sections
above and to the relevant non-functional requirements from
Chapter~3~\cite{owasp2021}.

\begin{table}[H]
  \centering
  \caption{OWASP Top Ten 2021 compliance mapping~\cite{owasp2021}}
  \label{tab:owasp}
  \begin{tabularx}{\linewidth}{llX}
    \toprule
    \textbf{OWASP ID} & \textbf{Category} & \textbf{Controls implemented} \\
    \midrule
    A01 & Broken Access Control
        & JWT authentication (Section~\ref{sec:auth}); ownership checks per resource; \texttt{requireAdmin} middleware; soft-delete with audit retention (Section~\ref{sec:authz}). \\
    A02 & Cryptographic Failures
        & bcrypt cost~12 for password hashing; HS256 JWT signing (RFC~7519~\cite{jwt2015}); HSTS via Helmet (Section~\ref{sec:headers}). \\
    A03 & Injection
        & Sequelize parameterised queries (Section~\ref{sec:validation}); \texttt{express-validator} input validation; Zod AI output schema validation. \\
    A05 & Security Misconfiguration
        & Helmet HTTP headers (Table~\ref{tab:helmet}); CORS allowlist (Section~\ref{sec:cors}); stack traces suppressed in production error handler. \\
    A07 & Auth Failures
        & Rate limiting 100 req/15 min/IP (Section~\ref{sec:rate_limit}); seven-day JWT expiry; email verification requirement; password reset TTL (Section~\ref{sec:auth}). \\
    A09 & Logging and Monitoring
        & Winston structured logging; \texttt{AuditLog} database table; login event recording (Section~\ref{sec:logging}). \\
    \bottomrule
  \end{tabularx}
\end{table}
